\documentclass[11pt,a4paper]{article}
\usepackage[margin=1in]{geometry}
\usepackage{amsmath,booktabs,hyperref,float}

\title{Iteration 016: Residual FIR Learning}
\author{Autoresearch Loop}
\date{April 8, 2026}

\begin{document}
\maketitle

\begin{abstract}
Learning a residual correction on top of the frozen closed-form solution
achieves $r=0.375$, not improving over the end-to-end FIR ($r=0.376$).
The residual signal is too weak for the FIR to reliably learn.
\end{abstract}

\section{Hypothesis}
Instead of fine-tuning the entire FIR from CF initialization, freeze the CF
spatial filter and learn only a residual FIR correction:
$\hat{y} = W^*x + \text{FIR}_\text{residual}(x)$. This decomposes the problem
into a known linear part and a learnable correction.

\section{Method}
Two-branch model: (1) frozen CF spatial filter, (2) trainable FIR filter (7 taps)
initialized to zero. Outputs are summed. Trained 150 epochs with channel dropout.

\section{Results}
\begin{table}[H]
\centering
\begin{tabular}{lrrr}
\toprule
Model & Mean $r$ & Std $r$ & SNR (dB) \\
\midrule
FIR + dropout (iter011) & 0.376 & 0.076 & 0.61 \\
CF + residual FIR & 0.375 & 0.077 & 0.61 \\
\bottomrule
\end{tabular}
\end{table}

\section{Analysis}
The residual approach performs identically to the baseline. The mapping is so
close to linear that the residual FIR learns nearly the same correction as
end-to-end fine-tuning from CF init. Residual decomposition adds complexity
without benefit when the base model is already near-optimal.

\section{Next}
Iteration 017: combined MSE+correlation loss with correlation-based validation.

\end{document}
